\subsection{Nagata 环}

定义 2. --- 称环 A 为 Nagata 环，如果它是 Noether 环，并且对于 A 的每个素理想 p，Noether 整环 A/p 都是日本环（第 1 目，定义 1）。

例. --- 1) 域上的每个有限生成代数都是 Nagata 环（第 1 目，例）。

\begin{enumerate}
\def\labelenumi{\arabic{enumi})}
\setcounter{enumi}{1}
\item
  每个完备局部 Noether 环都是 Nagata 环（第 2 目，定理 2）。
\item
  环 Z 是 Nagata 环（第 1 目，例以及命题 1 的推论）。
\item
  可以证明（习题 30），Nagata 环上的每个有限生成代数都是 Nagata 环。
\end{enumerate}

命题 4. --- 设 A 是一个 Nagata 环。

\begin{enumerate}
\def\labelenumi{\alph{enumi})}
\item
  每个有限 A-代数都是 Nagata 环。
\item
  对于 A 的每个乘性子集 S，环 \(\mathbf{S}^{-1}\mathbf{A}\) 都是 Nagata 环。
\item
  设 B 是一个有限 A-代数，\(\rho: A \to B\) 为典范同态。对于 B 的每个素理想 p，环 B/p 是日本环 \(A/\rho^{-1}(p)\) 上的一个有限代数，故为日本环（第 1 目，命题 2）。
\item
  设 q 是 \(S^{-1}A\) 的一个素理想；则存在 A 的素理想 p 使得 \(q = S^{-1}p\) 。环 ( \(S^{-1}A\) )/q 是日本环 A/p 的一个分式环，故为日本环（第 1 目，注记 2）。
\end{enumerate}

定理 3 (Zariski-Nagata). --- 设 A 是一个半局部 Noether 环。下列条件等价：

\begin{enumerate}
\def\labelenumi{(\roman{enumi})}
\item
  A 是 Nagata 环；
\item
  对于 A 的每个素理想 \(\mathfrak{p}\)，\(\kappa(\mathfrak{p})\) -代数 \(\kappa(\mathfrak{p}) \otimes_{\mathbf{A}} \hat{\mathbf{A}}\) 是可分的；
\item
  对于每个约化 A-代数 R，环 \(R \otimes_{A} \hat{A}\) 都是约化的。
\end{enumerate}

我们先证明条件 (ii) 与 (iii) 的等价性。蕴涵 (iii) \(\Rightarrow\) (ii) 是显然的；反之，设 A 满足条件 (ii)。于是，对于每个是域的 A-代数 K，环 K \(\otimes_{\mathbf{A}}\hat{\mathbf{A}}\) 都是约化的。现设 C 是一个约化的有限生成 A-代数；环 C 作为 Noether 环，同构于域的有限乘积 K \(_1\) \(\times\) \(\cdots\) \(\times\) K \(_n\) 的一个子环（IV，§ 2，第 5 目，命题 10）；由于 \(\hat{\mathbf{A}}\) 在 A 上平坦，环 C \(\otimes_{\mathbf{A}}\hat{\mathbf{A}}\) 同构于约化环 \(\prod_{i}(\mathrm{K}_{i}\otimes_{\mathbf{A}}\hat{\mathbf{A}})\) 的一个子环，故为约化的。最后，设 R 是任一约化 A-代数；则 R 是其有限生成子代数的过滤族 \((\mathrm{C}_{\alpha})\) 的并，且 \(R \otimes_{A} \hat{A}\) 是约化环的过滤族 \((\mathrm{C}_{\alpha} \otimes_{A} \hat{A})\) 的归纳极限，故为约化的。

我们来证明 (ii) 蕴涵 (i)。设 p 是 A 的一个素理想；环 A/p 的分式域 K 等同于 \(\kappa(p)\) ，而 K-代数 \(K \otimes_{A/p} (\widehat{A/p})\) 等同于 \(\kappa(p) \otimes_{A/p} \widehat{A}/p\widehat{A}\) ，从而等同于 \(\kappa(p) \otimes_{A} \widehat{A}\) 。若 \(\kappa(p) \otimes_{A} \widehat{A}\) 是一个 \(\kappa(p)\) -可分代数，则环 A/p 是日本环（第 2 目，命题 3）。

我们对 dim(A) 作归纳来证明蕴涵 (i) \(\Rightarrow\) (ii)。当 dim(A) = 0 时这是显然的，因为此时 A 是 Artin 环，从而是完备的。设 \(n\) 是一个整数 \(>0\) ；考虑如下的假设：

\((\mathbf{R}_n)\left\{
\begin{array}{ll}\text{对于每个局部诺特 Nagata 环 C，其维数满足 } < n\text{，且每个素}\\ \text{理想 } \mathfrak{r}\text{均为 C 的理想，则该环 }\kappa(\mathfrak{r})\otimes_{\mathrm{C}}\hat{\mathrm{C}}\text{是约化的。} \end{array}
\right.\)

设 A 是一个维数为 n 的半局部 Noether Nagata 环，p 是 A 的一个素理想，L 是域 \(\kappa(\mathfrak{p})\) 的一个有限扩张；在假设 (R \(_{n}\) ) 之下，只需证明环 \(L \otimes_{A} \hat{A}\) 是约化的。用 B 表示 A/p 在 L 中的整闭包；由于 A/p 是日本环，B 是一个有限 A-代数，因而是半局部 Nagata 环（命题 4）。设 m \(_{1}\) , \ldots, m \(_{r}\) 是 B 的极大理想；环 \(L \otimes_{A} \hat{A}\) 等同于 \(B \otimes_{A} \hat{A}\) 的一个分式环，而后者等同于诸局部环 B \(_{m_{i}}\) 的完备化之乘积（第 3 目，引理 1）。因此只需证明，对于 B 的每个极大理想 m，环 B \(_{m}\) 都是约化的（II，§ 2，第 6 目，命题 17）。环 B \(_{m}\) 是局部的、整闭的且是 Nagata 的（命题 4），并且有 dim(B \(_{m}\) ) ≤ dim(B) ≤ dim(A) = n（VIII，§ 1，第 3 目，命题 6 以及 § 2，第 3 目，定理 1）。改变记号，我们归结为在假设 (R \(_{n}\) ) 之下证明：对于每个维数 ≤ n 的局部 Noether 整闭 Nagata 环 A，环 \(\hat{A}\) 是约化的，也就是说（第 3 目，引理 2），\(\hat{A}_{p'}\) 对每个素理想 p′ ∈ Ass(\(\hat{A}\)) 都是约化的。由于当 dim(A) = 0 时这是直接的，我们可以假设 dim(A) > 0。于是设 x 是 m \(_{A}\) 的一个非零元素，q' 是 \(\hat{A}\) 的一个素理想，它在包含 \(x\hat{A} + p'\) 的素理想中极小；由于 \(\hat{A}_{p'}\) 等同于环 \(\hat{A}_{q'}\) 的一个分式环，只需证明后者是约化的（II，§ 2，第 6 目，命题 17）。由引理 3，理想 q' 与 \(\hat{A}\)-模 \(\hat{A}/x\hat{A}\) 相伴；由于 \(\hat{A}\) 在 A 上平坦，q' 在 A 中的逆像 q 与 A-模 A/xA 相伴（IV，§ 2，第 6 目，定理 2 的推论 1）。由于环 A 被假设为整闭的，这蕴涵 q 的高度为 1（VII，§ 1，第 6 目，命题 10），从而环 \(A_{q}\) 是一个离散赋值环（同前，第 3 目，定理 2 的推论以及第 6 目，定理 4）。由于 A/q 是一个维数 < n 的 Nagata 环，环 \(\kappa(q) \otimes_{A/q} \hat{A}/q\hat{A}\) 由假设 (R \(_{n}\) ) 是约化的。与它同构的环 \(\kappa(q) \otimes_{A} \hat{A}\) 是约化的，因而环 \(\kappa(q) \otimes_{A_{q}} \hat{A}_{q'}\) 作为它的一个分式环也是约化的。于是可以把第 3 节的引理 4 应用于典范同态 \(A_{q} \to \hat{A}_{q'}\)，从而得出环 \(\hat{A}_{q'}\) 是约化的，这就是我们想要证明的。定理 3 至此证毕。

推论 1. --- 局部且约化的 Nagata 环的完备化是约化的。

事实上，只需在定理 3 的 (iii) 中取 R = A。

推论 2 (Chevalley). --- 设 A 是域上的一个有限型约化代数，p 是 A 的一个素理想。局部环 A \(_{p}\) 的完备化是约化的。

由于 A 是约化的，局部环 \(A_{p}\) 是约化的；此外，A 是一个 Nagata 环（例 1），因而 \(A_{p}\) 是 Nagata 环（命题 4），而把推论 1 应用于环 \(A_{p}\) 即得推论 2。

推论 3. --- 设 k 是一个特征为 0 的域，A 是一个局部 Noether k-代数。A 为 Nagata 环的必要充分条件是：对于 A 的每个素理想 p，环 (A/p) 都是约化的。

事实上，由于诸域 \(\kappa(p)\) 的特征为 0，说代数 \(\kappa(p) \otimes_{A} \hat{A} = \kappa(p) \otimes_{A/p} (\widehat{A/p})\) 是约化的，与说它们是可分的是等价的（A，V，第 117 页，定理 1），这表明所述条件是充分的（定理 3，(ii) \(\Rightarrow\) (i)）；而且它也是必要的（定理 3，(i) \(\Rightarrow\) (iii)，取 \(R = A/p\) ）。

